\section{Введение}
\label{sec:Chapter0} \index{Chapter0}

%В этой части надо описать предметную область, задачу из которой вы будете
%решать, объяснить её актуальность (почему надо что-то делать сейчас?). Здесь же
%стоит ввести определения понятий, которые вам понадобятся в постановке задачи.

Большая часть современных криптографических протоколов базируются на сложности задач факторизации натурального числа и дискретного логарифмирования элемента в заданной абелевой группе.

Такими являются, например, широко известная криптосистема $\mathsf{RSA}$ ({\it Rivest–Shamir–Adleman})~\cite{RSA}, а также современные российские ГОСТы ({\it государственные стандарты}) для цифровой подписи документов~\cite{GOST_signature} и получения двумя сторонами общего секретного ключа~\cite{GOST_key_generation}.

При этом в $1994$\;-\,$1997$ годах Шор разработал алгоритм на квантовом компьютере, который решает обе эти задачи в $\N$ за полиномиальное от длины числа время~\cite{Shor}. В дальнейшем алгоритм улучшался и модифицировался, в том числе, в $2003$ году была опубликована и подробно описана его версия на эллиптических кривых~\cite{Ell_curves}.

Сейчас многими компаниями и государствами ведутся исследования по построению мощных стабильных квантовых компьютеров. В случае успеха оных все вышеупомянутые протоколы окажутся уязвимыми.

В связи с этим, начиная с $2016$\;-\:$2017$ годов NIST ({\it National Institute of Standards and Technology в США}) проводит конкурс по отбору лучших протоколов постквантовой ({\it устойчивой к атакам в том числе на квантовом компьютере}) криптографии~\cite{NIST_status}.

Немалая доля из них основана на решёточной криптографии вообще и $\mathsf{LWE}$-({\it Learning With Errors})-problem в частности~\cite{LWE}. Одним из главных преимуществ этих криптосистем является наличие $\mathsf{ACH}$ ({\it Average-Case}), а не $\mathsf{WCH}$ ({\it Worst-Case}) сложности ({\it Hardness})~\cite{ACH_Ajtai}.

Неформально говоря, это означает, что трудными для взлома являются случайные параметры протокола, а не худшие из всех возможных. В реальных криптосистемах данное свойство очень полезно, поскольку генерация рандомных бит при шифровании необходима.

К сожалению, протоколы, основанные $\mathsf{LWE}$-problem, хотя являются простыми и понятными, требуют довольно много вычислений. Поэтому с целью оптимизации в статье~\cite{RLWE} вводится $\mathsf{RLWE}$-({\it Ring LWE})-problem. Это модификация $\mathsf{LWE}$ в многочленах над $\Z_q$ c ускоренными вычислениями за счёт $\mathsf{NTT}$ ({\it Number Theoretic Transform})~\cite{NTT_applications}.

Более подробно все эти и некоторые дальнейшие концепции описаны в вышедшем в $2022$ году обзоре~\cite{Li2022ATI}. В моей же работе я буду использовать и улучшать результаты следующей статьи~\cite{RLWE_Channel}.

В ней процесс $\mathsf{RLWE}$-шифрования рассматривается, как передача данных по зашумлённому $\mathsf{RLWE}$-каналу. Исследуются свойства этого канала, описывается, как увеличить скорость передачи и при этом снизить $\DFR$ ({\it Decryption Failure Probability/Rate}) при расшифровке сообщения.

В частности, для уменьшения $\DFR$ при фиксированном $\mathsf{CFR}$ ({\it Coefficient Failure Rate}) используются $\mathsf{ECC}$ ({\it Error Correcting Codes})~\cite{Hamming}. Эта конструкция позволяет исправлять вплоть до $t$ покоординатных ошибок при передаче вектора-сообщения фиксированной длины $n$ (здесь $n, t$ --- некоторые выбираемые исследователем параметры).

При этом измерять кол-во произошедших ошибок можно по-разному. В моей работе я пробую разные способы (метрика Хэмминга~\cite{Hamming}, метрика Ли~\cite{Lee}) и вычисляю, как изменяется от этого скорость $\mathsf{RLWE}$-канала.

Как и в статье~\cite{RLWE_Channel}, я варьирую размер алфавита $Q$, которому принадлежат символы сообщения, и наблюдаю, как изменяются от этого оптимальные параметры канала.

Для этого я выбираю несколько конкретных криптосистем ($\LAC\,256$~\cite{LAC256}, $\NewHope\,1024$~\cite{NewHope1024}), целевое $\DFR$ и исследую их. В дальнейшем реализованный код и применённые методы можно обобщить также на другие криптосистемы с иными размерами алфавита и метриками.
